\section{\blue{结论}}
\label{sec:conclusion}

\blue{
预测模型正变得日益复杂，本文追问周期性传感器数据流是否需要这种复杂性。答案是否定的。周期序列的一个窗口携带三类信息——其各周期的形状、水平与振幅——它们构成一个低维结构：一个水平项加一个振幅缩放的形状，在七个基准数据集上捕获归一化窗口的 78\% 至 93\%。\sysplus 为每个部分赋予一个标记和一条小通路：相位标记——当每个周期以相同方式重塑时，它们停留在一个固定子空间中——通过一个带少量路由器与一个共享预测器的路由 Transformer 承载形状；水平标记即各周期均值，被映射到三个固定的预测时域轮廓上；振幅标记由水平计算，在可证明地保持每个未来周期均值不变的前提下缩放周期形状。统一的低维视角说明了三个标记为何缺一不可、每条通路为何取其形式：对固定权重，通过层归一化读取相位标记的预测模型，将其预测水平与振幅限制在两个任何附加深度都无法消除的包络内，而包络之外只是少数几个轮廓系数和一个由水平驱动的因子。从训练后的模型读出的包络，恰好在水平与振幅标记收益最大之处约束生效。

\sysplus 以其相位通路约 1.2K 参数量、水平与振幅至多再加约 2.1K 参数量，在七个基准数据集中的六个上取得了相对八个基线的最佳平均 MSE；在最大的基准数据集 Traffic 上，它将 PatchTST 和 Crossformer 的 FLOPs 削减约 99.99\%。由于一个共享模型服务每个传感器，这些节省会在网络每个传感器上、每次更新时重复兑现；而振幅标记承载的水平–振幅耦合在 Traffic 与 Electricity 这两个传感器网络上最强。更广泛地说，这些结果表明，预测模型所需的容量取决于其标记留下多少纠缠：有了覆盖预测所需内容的标记，小模型便已足够。

\fakeparagraph{局限与未来工作}
\sysplus 是通道独立的：它通过并行与一个共享模型来利用传感器网络的规模，而将其标记与显式空间建模相结合是自然的下一步。我们结果中关于能力的部分，是在季节性随机游走及其乘性变体下的存在性结果；而包络对每个参数取值都成立，并在训练后的模型上被测量。分析假设单一主导周期，以及为整数个周期的预测时域；多个相互作用的周期与无明显周期的序列，以及通路之间依赖输入的门控，都留待未来工作。
}
